%%%%%%%%%%%%%%%%%%%%%%%%%%%%%%%%%%%%%%%%%%%%%%%%%%%%%%%%%%%%%%%%%%%%%%%%
%% RESUMO                                                              %%
%%                                                                     %%
%% Forma (art. 40):                                                     %%
%%   I  título do pedido centralizado, separado do texto — impresso     %%
%%      automaticamente por \abrirresumo a partir de dados-do-pedido.tex;%%
%%   II conciso, preferencialmente entre 50 e 200 palavras, não         %%
%%      excedendo uma página.                                            %%
%%                                                                     %%
%% Conteúdo (art. 41):                                                   %%
%%   I   indicar o setor técnico e resumir o que foi exposto no         %%
%%       relatório, nas reivindicações e nos desenhos, representando    %%
%%       adequadamente as diferentes categorias de reivindicação;       %%
%%   II  permitir compreensão clara do problema técnico, da essência da %%
%%       solução e do uso principal;                                     %%
%%   III servir de instrumento de pré-seleção em buscas técnicas;       %%
%%   IV  NÃO fazer menção ao mérito ou ao valor do objeto — nada de     %%
%%       "solução inovadora", "grande vantagem", "excelente desempenho".%%
%%                                                                     %%
%% As diretrizes de exame (Resolução INPI/PR nº 124/2013, itens 5.01 e  %%
%% 5.02) lembram que muitas buscas consultam apenas título e resumo:    %%
%% coloque aqui as palavras-chave pelas quais o seu pedido deve ser      %%
%% encontrado.                                                           %%
%%                                                                     %%
%% Escreva em UM ÚNICO parágrafo, sem \pnum: o resumo não tem numeração %%
%% de parágrafos (o art. 26, II vale só para o relatório descritivo).   %%
%%                                                                     %%
%% `make verificar` conta as palavras e reclama se você sair da faixa   %%
%% de 50 a 200.                                                          %%
%%%%%%%%%%%%%%%%%%%%%%%%%%%%%%%%%%%%%%%%%%%%%%%%%%%%%%%%%%%%%%%%%%%%%%%%

A presente \nomedanatureza\ pertence ao setor técnico dos equipamentos de
irrigação agrícola e refere-se a um dispositivo de acionamento (1) para válvula
de irrigação, constituído por um corpo tubular (2) que aloja uma membrana
elástica (3) solidária a uma haste de acionamento (4) movimentada por um
atuador eletromecânico (6). A haste de acionamento (4) apresenta um ressalto de
encosto que atua sobre uma mola de retorno (5) pré-carregada e recebe um
assento de vedação de dupla face (9), que veda alternadamente o bocal de saída
e um orifício de alívio do corpo tubular (2). A pré-carga da mola de retorno
(5) estabelece duas posições extremas estáveis, de modo que o atuador
eletromecânico (6) é alimentado apenas durante a transição entre os estados de
abertura e de fechamento. Um sensor de pressão (7) e uma unidade de controle
(8) comandam a transição por pulsos de corrente, conforme a pressão da linha.
\ifinvencao
A \nomedanatureza\ compreende ainda o processo de controle de vazão executado
pela unidade de controle (8), com etapas de leitura, comparação e aplicação de
pulso.
\fi
